\documentclass[11pt]{article}
\usepackage[utf8]{inputenc}
\usepackage{amsmath}
\usepackage{amssymb}
\usepackage{booktabs}
\usepackage{geometry}
\geometry{margin=2.5cm}

\title{Simplified Fossen Model}
\author{}
\date{}

\begin{document}
\maketitle

\section*{Kinematic model}

\begin{equation}
\dot{\boldsymbol{\eta}} = \mathbf{J}(\boldsymbol{\eta})\,\mathbf{v}
\end{equation}

\section*{Kinetic model (simplified 3-DOF)}

\begin{equation}
\mathbf{M}\dot{\mathbf{v}} + \mathbf{C}(\mathbf{v})\mathbf{v} + \mathbf{D}(\mathbf{v})\mathbf{v} = \boldsymbol{\tau}
\end{equation}

Expanded form:

\begin{equation}
\begin{cases}
\dot{x} = u\cos\psi - v\sin\psi\\[4pt]
\dot{y} = u\sin\psi + v\cos\psi\\[4pt]
\dot{\psi} = r\\[4pt]
(m-X_{\dot{u}})\,\dot{u} - (m-Y_{\dot{v}})\,vr + X_u\,u = X_{P1}+X_{P2}\\[4pt]
(m-Y_{\dot{v}})\,\dot{v} + (m-X_{\dot{u}})\,ur + Y_v\,v = 0\\[4pt]
(I_z-N_{\dot{r}})\,\dot{r} - \big[(m-X_{\dot{u}})-(m-Y_{\dot{v}})\big]uv + N_r\,r = (X_{P1}-X_{P2})\cdot d_P
\end{cases}
\end{equation}

\section*{The 6 model parameters}

\begin{table}[h]
\centering
\begin{tabular}{clll}
\toprule
\# & Parameter & Description & Value \\
\midrule
1 & $m_{11}=m-X_{\dot{u}}$ & Effective mass in surge (mass + added mass) & $50.05$ kg \\
2 & $m_{22}=m-Y_{\dot{v}}$ & Effective mass in sway & $84.36$ kg \\
3 & $m_{33}=I_z-N_{\dot{r}}$ & Effective inertia in yaw & $17.21$ kg$\cdot$m$^2$ \\
4 & $X_u$ & Linear drag coefficient in surge & $151.57$ N$\cdot$s/m \\
5 & $Y_v$ & Linear drag coefficient in sway & $132.5$ N$\cdot$s/m \\
6 & $N_r$ & Linear drag coefficient in yaw & $34.56$ N$\cdot$m$\cdot$s \\
\bottomrule
\end{tabular}
\end{table}

\section*{Distance between thrusters}

\begin{table}[h]
\centering
\begin{tabular}{cll}
\toprule
Parameter & Description & Value \\
\midrule
$d_P$ & Transverse distance from the USV centerline to each thruster & $0.26$ m \\
$B$ & Total distance between the two thrusters ($=2d_P$) & $0.52$ m \\
\bottomrule
\end{tabular}
\end{table}

\section*{Thrust allocation matrix}

Since the USV has only two thrusters, both aligned with the surge direction ($\alpha_1=\alpha_2=0$), the control vector reduces to two useful components ($\tau_v=0$ at all times, as there is no lateral component):

\begin{equation}
  \boldsymbol{\tau} = \mathbf{B}\,\mathbf{T}
\end{equation}

where

\begin{equation}
  \boldsymbol{\tau} = \begin{bmatrix}\tau_u\\ \tau_r\end{bmatrix}, \qquad
  \mathbf{T} = \begin{bmatrix}T_1\\ T_2\end{bmatrix}, \qquad
  \mathbf{B} = \begin{bmatrix}1 & 1\\ d_P & -d_P\end{bmatrix}
\end{equation}

That is:

\begin{equation}
  \begin{bmatrix}\tau_u\\ \tau_r\end{bmatrix}
  =
  \begin{bmatrix}1 & 1\\ d_P & -d_P\end{bmatrix}
  \begin{bmatrix}T_1\\ T_2\end{bmatrix}
\end{equation}

Here $T_1$ and $T_2$ are the port and starboard thruster forces, respectively, and $d_P = 0.26$ m is the transverse distance from each thruster to the USV centerline. Since the number of control inputs equals the number of useful outputs, $\mathbf{B}$ is square (2$\times$2) and the system is not overactuated.

\section*{Thruster thrust model (cmd--thrust curve)}

The actual thruster thrust model is a generalized logistic function, with a dead zone for $|cmd|\leq 0.01$:

\begin{equation}
T =
\begin{cases}
0, & |cmd| \leq 0.01 \\[8pt]
A_{pos} + \dfrac{K_{pos} - A_{pos}}{\left(C_{pos} + \exp\left(-B_{pos} (cmd - M_{pos})\right)\right)^{\frac{1}{v_{pos}}}}, & cmd > 0.01 \\[12pt]
A_{neg} + \dfrac{K_{neg} - A_{neg}}{\left(C_{neg} + \exp\left(-B_{neg} (cmd - M_{neg})\right)\right)^{\frac{1}{v_{neg}}}}, & cmd < -0.01
\end{cases}
\label{equa:cmd2thrust}
\end{equation}

\begin{table}[h]
\centering
\caption{Optimized parameters for the thrust model}
\label{tab:t200}
\begin{tabular}{|c|c||c|c|}
\hline
\textbf{Parameter} & \textbf{Value} & \textbf{Parameter} & \textbf{Value} \\
\hline
$A_{pos}$ & 0.000001 & $A_{neg}$ & -31.4990 \\
$K_{pos}$ & 40.0209  & $K_{neg}$ & -0.00001 \\
$B_{pos}$ & 2.6249   & $B_{neg}$ & 3.6986 \\
$v_{pos}$ & 0.1615   & $v_{neg}$ & 0.3264 \\
$C_{pos}$ & 0.9432   & $C_{neg}$ & 0.9713 \\
$M_{pos}$ & 0.00001  & $M_{neg}$ & -1.0000 \\
\texttt{maxForceFwd} & 36.3827 N & \texttt{maxForceRev} & -28.4393 N \\
\hline
\end{tabular}
\end{table}

\subsection*{5th-degree polynomial approximation}

To simplify implementation in the control model, the logistic curve was approximated with a 5th-degree polynomial over the range $cmd \in [-1,1]$, obtained via least-squares fitting:

\begin{equation}
T(cmd) \approx -26.557\,cmd^5 + 0.415\,cmd^4 + 46.293\,cmd^3 + 3.732\,cmd^2 + 11.674\,cmd + 0.243
\end{equation}


\end{document}